\chapter{Ocular Histoplasmosis Syndrome}
\index{Ocular Histoplasmosis Syndrome}
\label{sec:Ocular Histoplasmosis Syndrome}

\section{Overview}
\begin{itemize}[noitemsep]
  \item Ocular histoplasmosis syndrome (OHS) is an eye condition caused by the fungus Histoplasma capsulatum, which infects the eye years after a lung infection
  \item It causes punched-out scars in the retina and choroid and can lead to choroidal neovascularization with central vision loss
  \item It is most common in the Ohio and Mississippi River valleys of the United States, where the fungus is endemic
\end{itemize}
\section{Causes}
This condition happens because of an earlier, often silent respiratory infection with Histoplasma capsulatum. The fungus reaches the eye through the bloodstream, where it establishes dormant lesions in the choroid. Years later, inflammation at these sites can trigger the growth of new, fragile blood vessels (choroidal neovascularization) that leak and scar the macula.
\section{Risk Factors}
\begin{itemize}[noitemsep]
  \item Living or having lived in endemic regions (Ohio/Mississippi River valleys)
  \item Exposure to soil contaminated with bird or bat droppings
  \item Previous pulmonary histoplasmosis (often asymptomatic)
  \item Age over 40 years
  \item Preexisting ocular histoplasmosis scars
\end{itemize}
\section{Signs and Symptoms}
\subsection{Common symptoms}
\begin{itemize}[noitemsep]
  \item Many people have no symptoms and lesions are found on routine examination
  \item Blurred or distorted central vision when the macula is involved
  \item Central blind spot (scotoma) or metamorphopsia
  \item In peripheral lesions, vision may be normal
\end{itemize}
\subsection{Emergency warning signs}
\begin{itemize}[noitemsep]
  \item Sudden worsening of central vision or distortion requires urgent evaluation, as early treatment of neovascularization preserves vision
\end{itemize}
\section{Progression and Stages}
Histoplasmosis infection is usually acquired in childhood or young adulthood with no symptoms. Ocular lesions (punched-out chorioretinal scars, peripapillary atrophy, and linear streaks) may be present for years before any visual symptoms. The vision-threatening stage occurs when choroidal neovascularization develops, usually in the fourth to sixth decade, causing fluid, hemorrhage, and scarring in the macula.
\section{Complications}
\begin{itemize}[noitemsep]
  \item Choroidal neovascularization with macular hemorrhage and scarring
  \item Permanent central vision loss
  \item Recurrent neovascularization after treatment
  \item Scarring and atrophy of the retina
  \item Reduced quality of life from impaired central vision
\end{itemize}
\section{Diagnosis}
\begin{itemize}[noitemsep]
  \item Dilated fundus examination showing the characteristic punched-out scars
  \item Optical coherence tomography (OCT) to detect fluid or neovascularization
  \item Fluorescein angiography to identify leaking new vessels
  \item Indocyanine green angiography in selected cases
  \item Histoplasmosis skin or antibody testing confirms prior exposure but is not required for diagnosis
\end{itemize}
\section{Prevention}
\begin{itemize}[noitemsep]
  \item Avoid exposure to disturbed soil in endemic areas, especially with bird or bat droppings
  \item Use protective measures when working in high-risk environments (dust masks)
  \item Regular eye examinations for people with known histoplasmosis scars
  \item Early self-monitoring with an Amsler grid to detect distortion
\end{itemize}
\section{Treatment Options}
\begin{itemize}[noitemsep]
  \item Observation for inactive lesions without neovascularization
  \item Intravitreal anti-VEGF injections (ranibizumab, aflibercept, bevacizumab) as first-line therapy for active neovascularization
  \item Photodynamic therapy (PDT) as an adjunct in selected cases
  \item Laser photocoagulation for well-defined extrafoveal lesions (less commonly used today)
  \item Systemic or local antifungal therapy is generally not helpful once ocular lesions are established
  \item Regular monitoring of the fellow eye
\end{itemize}
\section{Living with It}
\begin{itemize}[noitemsep]
  \item Use an Amsler grid at home to monitor for early distortion or scotomas
  \item Attend regular follow-up for both eyes
  \item Use low-vision aids if central vision is impaired
  \item Report any new visual symptoms promptly
\end{itemize}
\section{Outlook}
With modern anti-VEGF therapy, many patients with neovascular OHS retain or regain useful central vision, though repeated injections are often needed. Peripheral lesions have an excellent prognosis. Untreated central neovascularization typically causes significant permanent vision loss.
